%%%%%%%%%%%%%%%%%%%%%%%%%%%%%%%%%%%%%%%%%%%%%%%%%%%%%%%
% NOTAS TEORIA DE TRANSPORTE OPTIMO
%
% Contains ?? figures. 
% File started: ??
% First circulated version: ??
% Last version:
% Last modification: -
%%%%%%%%%%%%%%%%%%%%%%%%%%%%%%%%%%%%%%%%%%%%%%%%%%%%%%%

\documentclass[11pt,a4paper]{amsart}
\usepackage{xr-hyper}
\newcounter{chapter}

\usepackage{amssymb}
\usepackage{hyperref}
\usepackage{listings}
\usepackage{mathtools}
\usepackage{enumitem}

% Fuzz -------------------------------------------------------------------
\hfuzz2pt % Don't bother to report over-full boxes if over-edge is < 2pt
% Line spacing -----------------------------------------------------------

%%%%%%%%%fin pegada%%%%%%%%%%%%%%%%%%5

\usepackage[T1]{fontenc}
\usepackage[spanish]{babel}
\usepackage{graphicx}
\usepackage{tikz}
\usetikzlibrary{babel}

\setlength{\textwidth}{15.5cm} 
\setlength{\textheight}{24cm}
\setlength{\oddsidemargin}{0cm}
\setlength{\evensidemargin}{0cm}


\parskip 3pt

%%
%% OPERADORES
%%
\renewcommand{\Re}{\operatorname{Re}}
\newcommand{\tr}{\operatorname{tr}}
\newcommand{\id}{\operatorname{id}}
\def\div{\operatorname {\mathrm{div}}}
\def\Dom{\operatorname {\mathrm{Dom}}}
\def\argmin{\operatorname {\mathrm{arg\, min}}}
\def\argmax{\operatorname {\mathrm{arg\, max}}}
\def\gen{\operatorname {\mathrm{gen}}}
\def\Im{\operatorname {\mathrm{Im}}}
\def\Nu{\operatorname {\mathrm{Nu}}}
\def\sop{\operatorname {\mathrm{sop}}}
\def\diam{\operatorname {\mathrm{diam}}}
\def\dist{\operatorname {\mathrm{dist}}}
\def\var{\operatorname {\mathrm{Var}}}
\def\codim{\operatorname {\mathrm{codim}}}
\def\pint{\operatorname {--\!\!\!\!\!\int\!\!\!\!\!--}}

%%
%% EPSILON
%%
\newcommand{\ve}{\varepsilon}

%%
%% LETRAS MATHBB
%%
\newcommand{\R}{{\mathbb R}}
\newcommand{\Z}{{\mathbb Z}}
\newcommand{\C}{{\mathbb C}}
\newcommand{\Q}{{\mathbb Q}}
\newcommand{\N}{{\mathbb N}}
\newcommand{\E}{{\mathbb E}}
\newcommand{\V}{{\mathbb V}}
\renewcommand{\P}{{\mathbb P}}
\newcommand{\II}{{\mathbb I}}
\newcommand{\OO}{{\mathbb O}}
\newcommand{\Sb}{{\mathbb S}}
\newcommand{\Tb}{{\mathbb T}}

%%
%% LETRAS CALIFGRAFICAS
%%
\newcommand{\F}{{\mathcal F}}
\newcommand{\Hc}{{\mathcal H}}
\newcommand{\Tc}{{\mathcal T}}
\newcommand{\Ec}{{\mathcal E}}
\newcommand{\Vc}{{\mathcal V}}
\newcommand{\A}{{\mathcal A}}
\newcommand{\J}{{\mathcal J}}
\newcommand{\G}{{\mathcal G}}
\newcommand{\U}{{\mathcal U}}
\newcommand{\B}{{\mathcal B}}
\newcommand{\Sw}{{\mathcal S}}
\newcommand{\D}{{\mathcal D}}
\newcommand{\M}{{\mathcal M}}
\newcommand{\LL}{{\mathcal L}}

%%
%% LETRAS MATHBF
%%
\newcommand{\ab}{{\mathbf a}}
\newcommand{\bb}{{\mathbf b}}
\newcommand{\cc}{{\mathbf c}}
\newcommand{\ee}{{\mathbf e}}
\newcommand{\ff}{{\mathbf f}}
\newcommand{\gb}{{\mathbf g}}
\newcommand{\hh}{{\mathbf h}}
\newcommand{\n}{{\mathbf n}}
\newcommand{\mm}{{\mathbf m}}
\newcommand{\pp}{{\mathbf p}}
\newcommand{\qq}{{\mathbf q}}
\newcommand{\sbf}{{\mathbf s}}
\newcommand{\tb}{{\mathbf t}}
\newcommand{\uu}{{\mathbf u}}
\newcommand{\vv}{{\mathbf v}}
\newcommand{\ww}{{\mathbf w}}
\newcommand{\xx}{{\mathbf x}}
\newcommand{\yy}{{\mathbf y}}
\newcommand{\zz}{{\mathbf z}}
\newcommand{\I}{{\mathbf 1}}
\newcommand{\Pb}{{\mathbf P}}
\newcommand{\XX}{{\mathbf X}}
\newcommand{\YY}{{\mathbf Y}}

%%
%% LETRAS MATHFRAK
%%

\newcommand{\Gk}{{\mathfrak G}}


%%
%% TEOREMAS
%%
\newtheorem{teo}{Teorema}[section]
\newtheorem{lema}[teo]{Lema}
\newtheorem{prop}[teo]{Proposición}
\newtheorem{cor}[teo]{Corolario}

%%
%% REMARK
%%
\theoremstyle{remark}
\newtheorem{obs}[teo]{Observación}

%%
%% DEFINICION
%%
\theoremstyle{definition}
\newtheorem{defi}[teo]{Definición}
\newtheorem{ejem}[teo]{Ejemplo}
\newtheorem{ejer}[teo]{Ejercicio}
\newtheorem{nota}[teo]{Notación}
\newtheorem{ejercicio}{Ejercicio}[chapter]


%%
%% NUMERACION
%%
\numberwithin{subsection}{section}
\numberwithin{equation}{section}

%%
%% TITULO
%%


\externaldocument[N-]{../notas/Notas-TO}
\newcommand{\cuadernolink}[2]{\noindent\textbf{Cuaderno:} \emph{#1} (\href{https://colab.research.google.com/github/jfbonder/transporte-optimo/blob/main/cuadernos/#2}{abrir en Colab}).\par\medskip}
\pagestyle{plain}

\begin{document}
\renewcommand{\thechapter}{4}
\begin{center}
{\large\sc Teoría de Transporte Óptimo}\\[2pt]
{\small 2do cuatrimestre 2026 --- Julián Fernández Bonder}\\[10pt]
{\LARGE\bf Práctica 4}\\[4pt]
{\large Dualidad}
\end{center}
\bigskip
\noindent{\small Los ejercicios son los de la sección de ejercicios del Capítulo~4 de las notas del curso, con la misma numeración. Las referencias a teoremas, proposiciones y secciones remiten a las notas (\url{https://github.com/jfbonder/transporte-optimo}).}
\medskip


\begin{ejercicio}[Desigualdades ``casi en todo punto'' y planes]\label{ejer.ctp.planes}
Sean $\pi\in\Pi(\mu,\nu)$ y $c$, $\phi$, $\psi$ funciones medibles.
\begin{enumerate}[label=(\alph*)]
\item Supongamos que existen $N\subset X$ y $M\subset Y$ con
$\mu(N)=\nu(M)=0$ tales que $\phi(x)+\psi(y)\le c(x,y)$ para todo $x\notin N$
y todo $y\notin M$. Probar que $\phi(x)+\psi(y)\le c(x,y)$ para
$\pi$-casi todo $(x,y)$.
\item Mostrar que la conclusión es \emph{falsa} si sólo se supone que la
desigualdad vale $(\mu\otimes\nu)$-c.t.p.: tomar $X=Y=[0,1]$,
$\mu=\nu=\LL^1$, $c=\I_{\{x\ne y\}}$, $\phi=\psi=\tfrac12$ y
$\pi=(\id,\id)_\#\mu$.
\item Explicar por qué (b) obliga a definir la clase admisible
$\mathcal A_c$ con la desigualdad en \emph{todo} punto, y qué fallaría en la
dualidad débil (Proposición~\ref{N-prop.dualidad.debil}) con una definición
``c.t.p.''.
\end{enumerate}
\end{ejercicio}

\begin{ejercicio}[Dualidad a mano]\label{ejer.dualidad.a.mano}
Sean $\mu=\LL^1\lfloor_{[0,1]}$ y $\nu=\LL^1\lfloor_{[1,2]}$.
\begin{enumerate}[label=(\alph*)]
\item Para $c(x,y)=|x-y|$, probar que $(\varphi,\psi)=(-x,y)$ es admisible
y que el plan $\pi=(\id,T)_\#\mu$ con $T(x)=x+1$ tiene costo igual a
$\D(\varphi,\psi)$. Concluir que ambos son óptimos y que el costo óptimo es
$1$. ¿Hay otros planes óptimos?
\item Para $c(x,y)=|x-y|^2$, hallar un par admisible $(\varphi,\psi)$ que
sature en el gráfico de $T$ (buscar $\psi$ afín) y concluir que $T$ es
óptimo también para el costo cuadrático. ¿Es único el plan óptimo en este
caso?
\item Interpretar $\varphi$ y $\psi$ en términos del problema del
transportista (precios de carga y descarga) de este capítulo. Observar que
en (a) \emph{todo} plan de $\Pi(\mu,\nu)$ es óptimo; explicar por qué.
\end{enumerate}
\end{ejercicio}

\begin{ejercicio}[Funciones semicontinuas y medidas signadas]\label{ejer.sci.signadas}
Sea $X$ un espacio métrico compacto y $f\colon X\to\R\cup\{+\infty\}$
semicontinua inferiormente y acotada inferiormente. Para una medida de
Borel signada finita $\mu$ probar que
\[
\sup\Bigl\{\int_Xg\,d\mu:\ g\in C_b(X),\ g\le f\Bigr\}
=\begin{cases}\displaystyle\int_Xf\,d\mu,&\text{si }\mu\ge0,\\[1ex]
+\infty,&\text{en otro caso.}\end{cases}
\]
\emph{Sugerencia:} para la primera alternativa usar el
Ejercicio~\ref{N-ejer.sci.integral}; para la segunda, la descomposición de
Hahn de $\mu$ y funciones continuas que aproximen a
$-k\,\I_{A}$ sobre un conjunto $A$ con $\mu(A)<0$. Este cálculo es el que
identifica la transformada de Legendre--Fenchel de la restricción
$\varphi\oplus\psi\le c$ en la demostración del Teorema~\ref{N-teo.dualidad}.
\end{ejercicio}

\begin{ejercicio}[$\bigstar$ Dualidad para costos semicontinuos inferiores]\label{ejer.dualidad.sci}
En la Sección~\ref{N-sec.demo.dualidad} se demuestra el Teorema~\ref{N-teo.dualidad}
para $X$, $Y$ compactos y $c$ continua. Extenderlo a $c\colon X\times Y\to[0,+\infty]$
semicontinua inferiormente, con $X$, $Y$ compactos, siguiendo estos pasos.
\begin{enumerate}[label=(\alph*)]
\item Por el Ejercicio~\ref{N-ejer.moreau.yosida}, existen $c_k\in C(X\times Y)$
con $c_k\nearrow c$. Sean $\pi_k$ planes óptimos para $c_k$. Probar que toda
subsucesión de $(\pi_k)$ tiene una subsucesión convergente a un plan $\pi$ y
que, para todo $k_0$ fijo,
$\int c_{k_0}\,d\pi\le\lim_k\min_{\Pi}\int c_k\,d\gamma$. Deducir, con
convergencia monótona, que $\int c\,d\pi\le\lim_k\min_\Pi\int c_k\,d\gamma$,
y por lo tanto que $\min_\Pi\int c_k\,d\gamma\nearrow\min_\Pi\int c\,d\gamma$.
\item Probar que $\mathcal A_{c_k}\subset\mathcal A_c$ y concluir
\[
\sup_{\mathcal A_c}\D\ \ge\ \sup_k\sup_{\mathcal A_{c_k}}\D
=\sup_k\min_\Pi\int c_k\,d\gamma=\min_\Pi\int c\,d\gamma\ \ge\ \sup_{\mathcal A_c}\D,
\]
donde la última desigualdad es la dualidad débil.
\item ¿Qué cambia si $X$, $Y$ son sólo completos y separables? Consultar
\cite[Thm.~1.3]{Villani2003} o \cite[Thm.~1.39]{Santambrogio}.
\end{enumerate}
\end{ejercicio}

\bibliographystyle{plain}
\bibliography{../notas/biblio}

\end{document}
